\paragraph{bin-fold 群胚代数的有序 \texorpdfstring{$K_0$}{K0}、稳定刚性与自同构阈值}
在 bin-fold 纤维群胚代数的 Wedderburn 信息完备性、Morita--\(K\) 理论、PI 阈值与 window-\(6\) 的异常资源分叉已经建立之后，还可把同一退化谱直方图继续压缩为如下六条对象层结论：带单位元类的有序 \(K_0\) 完全分类、稳定同构刚性、PI 度数的黄金下界、自同构群的半直积分解、中心块数与异常比特数的严格对齐，以及首次非交换阈值的精确定位。

\begin{theorem}[bin-fold 纤维群胚代数的有序 \texorpdfstring{$K_0$}{K0} 完全分类]\label{thm:xi-time-part60acb-binfold-ordered-k0-complete-classification}
沿用定义 \ref{def:fold-bin-fiber-groupoid} 与命题 \ref{prop:fold-bin-groupoid-algebra-wedderburn} 的记号，记
\[
A_m:=\CC[\mathcal R_m^{\mathrm{bin}}],
\qquad
\mathcal R_m^{\mathrm{bin}}
=
\bigsqcup_{w\in X_m}(F_w\times F_w),
\qquad
d_w:=|F_w|=d_m^{\mathrm{bin}}(w).
\]
则在由 Wedderburn 块给出的自然坐标下，有
\[
\boxed{\;
\bigl(K_0(A_m),K_0(A_m)^+,[1_{A_m}]\bigr)
\cong
\bigl(\ZZ^{X_m},\NN^{X_m},(d_w)_{w\in X_m}\bigr).\;}
\]
特别地，带 order unit 的有序 \(K_0\) 唯一恢复纤维多重度多重集
\[
\{d_w:w\in X_m\},
\]
反过来，该多重度多重集也唯一决定 \(A_m\) 的复代数同构类。
\end{theorem}
\begin{proof}
由命题 \ref{prop:fold-bin-groupoid-algebra-wedderburn}，
\[
A_m\cong \bigoplus_{w\in X_m}M_{d_w}(\CC).
\]
因此
\[
K_0(A_m)\cong \bigoplus_{w\in X_m}K_0\!\bigl(M_{d_w}(\CC)\bigr).
\]
对每个矩阵块 \(M_{d_w}(\CC)\)，以秩一投影类作为 \(K_0\)-生成元，则
\[
K_0\!\bigl(M_{d_w}(\CC)\bigr)\cong \ZZ,
\qquad
K_0\!\bigl(M_{d_w}(\CC)\bigr)^+\cong \NN,
\qquad
[1_{M_{d_w}(\CC)}]=d_w.
\]
逐块直和便得到
\[
K_0(A_m)\cong \ZZ^{X_m},
\qquad
K_0(A_m)^+\cong \NN^{X_m},
\qquad
[1_{A_m}]=(d_w)_{w\in X_m}.
\]

由于 \(\NN^{X_m}\) 的极端射线恰对应于各个 Wedderburn 块，故在不固定块顺序时，order unit 的坐标仍唯一恢复多重度多重集 \(\{d_w\}_{w\in X_m}\) 仅差一个置换。反向方向由定理 \ref{thm:xi-time-part60aaa-wedderburn-information-completeness} 立即给出：纤维直方图，等价地多重度多重集，唯一决定 \(A_m\) 的半单代数同构型。证毕。
\end{proof}

\begin{corollary}[Morita 压缩与稳定刚性]\label{cor:xi-time-part60acb-binfold-morita-stable-rigidity}
设 \(\mathcal K\) 为可分无限维 Hilbert 空间上的紧算子代数。则对每个 \(m\ge 1\)，有
\[
\boxed{\;
A_m\ \text{与}\ \CC^{F_{m+2}}\ \text{Morita 等价}.\;}
\]
并且对任意 \(m,m'\ge 1\)，成立等价
\[
\boxed{\;
A_m\otimes\mathcal K\cong A_{m'}\otimes\mathcal K
\iff
F_{m+2}=F_{m'+2}
\iff
m=m'.\;}
\]
\end{corollary}
\begin{proof}
第一条正是定理 \ref{thm:xi-time-part60aca-morita-type-forgets-block-sizes} 的内容。

另一方面，由 Wedderburn 分解，
\[
A_m\otimes\mathcal K
\cong
\bigoplus_{w\in X_m}\bigl(M_{d_w}(\CC)\otimes\mathcal K\bigr)
\cong
\bigoplus_{w\in X_m}\mathcal K
\cong
\mathcal K^{\oplus F_{m+2}},
\]
因为 \(M_d(\CC)\otimes\mathcal K\cong \mathcal K\) 对每个 \(d\ge 1\) 都成立。
故若
\[
A_m\otimes\mathcal K\cong A_{m'}\otimes\mathcal K,
\]
则由稳定性
\[
K_0(A_m\otimes\mathcal K)\cong K_0(A_m),
\qquad
K_0(A_{m'}\otimes\mathcal K)\cong K_0(A_{m'}),
\]
再结合定理 \ref{thm:xi-time-part60acb-binfold-ordered-k0-complete-classification}，得到
\[
\ZZ^{F_{m+2}}\cong K_0(A_m)\cong K_0(A_{m'})\cong \ZZ^{F_{m'+2}},
\]
从而
\[
F_{m+2}=F_{m'+2}.
\]
反之，若两者相等，则两边都同构于同一个
\(
\mathcal K^{\oplus F_{m+2}}
\)。

最后，由 Fibonacci 数列严格递增，\(F_{m+2}=F_{m'+2}\) 等价于 \(m=m'\)。证毕。
\end{proof}

\begin{theorem}[PI 度数的纤维极值律与黄金下界]\label{thm:xi-time-part60acb-binfold-pideg-golden-lower-bound}
沿用记号
\[
A_m\cong \bigoplus_{w\in X_m}M_{d_w}(\CC),
\qquad
M_m:=\max_{w\in X_m}d_w,
\qquad
\overline d_m:=\frac{2^m}{F_{m+2}}.
\]
则
\[
\boxed{\;
\operatorname{PIdeg}(A_m)=M_m.\;}
\]
等价地，\(A_m\) 的最小标准多项式恒等式阶数恰为
\[
\boxed{\;
2M_m.\;}
\]
进一步，若 \(C_\varphi>0\) 为定理 \ref{thm:fold-critical-resonance-constant} 的黄金共振常数，则
\[
\boxed{\;
\liminf_{m\to\infty}
\frac{\operatorname{PIdeg}(A_m)}{\overline d_m}
\ge
1+C_\varphi.\;}
\]
\end{theorem}
\begin{proof}
由定理 \ref{thm:xi-time-part60aca-pideg-maxfiber-threshold}，已知
\[
\operatorname{PIdeg}(A_m)=M_m,
\qquad
A_m\models s_{2M_m},
\qquad
A_m\not\models s_{2M_m-2},
\]
故前两条断言成立。

再由定理 \ref{thm:xi-time-part9m-maxfiber-uniformity-gap-chain}，
\[
\liminf_{m\to\infty}\frac{M_m}{\overline d_m}\ge 1+C_\varphi.
\]
把 \(\operatorname{PIdeg}(A_m)=M_m\) 代入即得
\[
\liminf_{m\to\infty}
\frac{\operatorname{PIdeg}(A_m)}{\overline d_m}
\ge
1+C_\varphi.
\]
证毕。
\end{proof}

\begin{theorem}[纤维群胚代数自同构群的半直积分解]\label{thm:xi-time-part60acb-binfold-algebra-automorphism-semidirect}
对每个 \(d\ge 1\) 记
\[
n_d(m):=\#\{w\in X_m:\ d_w=d\}.
\]
则
\[
\boxed{\;
\operatorname{Aut}_{\CC\text{-alg}}(A_m)
\cong
\left(\prod_{\substack{d\ge 1\\ n_d(m)\neq 0}}\operatorname{PGL}_d(\CC)^{\,n_d(m)}\right)
\rtimes
\left(\prod_{\substack{d\ge 1\\ n_d(m)\neq 0}}\mathfrak S_{n_d(m)}\right),\;}
\]
其中右端半直积作用由各 \(\mathfrak S_{n_d(m)}\) 对同尺寸 \(\operatorname{PGL}_d(\CC)\)-因子的置换给出。

特别地，在 \(m=6\) 时，由直方图
\[
2:8,\qquad 3:4,\qquad 4:9,
\]
有
\[
\boxed{\;
\operatorname{Aut}_{\CC\text{-alg}}(A_6)
\cong
\bigl(\operatorname{PGL}_2(\CC)^8\times \operatorname{PGL}_3(\CC)^4\times \operatorname{PGL}_4(\CC)^9\bigr)
\rtimes
\bigl(\mathfrak S_8\times \mathfrak S_4\times \mathfrak S_9\bigr).\;}
\]
并且
\[
\boxed{\;
\pi_0\!\bigl(\operatorname{Aut}_{\CC\text{-alg}}(A_6)\bigr)
\cong
\mathfrak S_8\times \mathfrak S_4\times \mathfrak S_9.\;}
\]
\end{theorem}
\begin{proof}
由定理 \ref{thm:xi-time-part60aaa-wedderburn-information-completeness}，
\[
A_m
\cong
\bigoplus_{\substack{d\ge 1\\ n_d(m)\neq 0}}M_d(\CC)^{\oplus n_d(m)}.
\]
由于 \(M_d(\CC)\not\cong M_e(\CC)\)（\(d\neq e\)）作为简单复代数彼此不同构，任一代数自同构都不能在不同尺寸的块之间混合；故
\[
\operatorname{Aut}_{\CC\text{-alg}}(A_m)
\cong
\prod_{\substack{d\ge 1\\ n_d(m)\neq 0}}
\operatorname{Aut}_{\CC\text{-alg}}\!\bigl(M_d(\CC)^{\oplus n_d(m)}\bigr).
\]

现固定一个 \(d\)，记 \(n:=n_d(m)\) 且
\[
B_d:=M_d(\CC)^{\oplus n}.
\]
任一 \(\CC\)-代数自同构 \(\Phi\) 必将 \(B_d\) 的极小中心幂等元两两置换，从而给出一个排列
\[
\sigma\in \mathfrak S_n.
\]
对每个分量，\(\Phi\) 在 \(M_d(\CC)\) 上诱导一个复代数自同构；由 Skolem--Noether 定理，每个此类自同构都由某个
\[
g_i\in \operatorname{PGL}_d(\CC)
\]
给出。于是
\[
\operatorname{Aut}_{\CC\text{-alg}}(B_d)
\cong
\operatorname{PGL}_d(\CC)^n\rtimes \mathfrak S_n,
\]
其中 \(\mathfrak S_n\) 通过置换 \(n\) 个等尺寸分量作用。
对全部 \(d\) 同时取直积，即得所示总分解。

当 \(m=6\) 时，推论 \ref{cor:fold-bin-m6-fiber-hist} 给出
\[
n_2(6)=8,\qquad n_3(6)=4,\qquad n_4(6)=9,
\]
代入即得显式公式。

最后，\(\operatorname{PGL}_d(\CC)\) 在复拓扑下连通，而 \(\mathfrak S_n\) 离散，故
\[
\pi_0\!\bigl(\operatorname{PGL}_d(\CC)^n\rtimes \mathfrak S_n\bigr)\cong \mathfrak S_n.
\]
对 \(d=2,3,4\) 取直积，便得到
\[
\pi_0\!\bigl(\operatorname{Aut}_{\CC\text{-alg}}(A_6)\bigr)
\cong
\mathfrak S_8\times \mathfrak S_4\times \mathfrak S_9.
\]
证毕。
\end{proof}

\begin{theorem}[window-\texorpdfstring{$6$}{6} 的中心块数与异常比特数严格对齐]\label{thm:xi-time-part60acb-window6-center-anomaly-bit-alignment}
记
\[
A_6:=\CC[\mathcal R_6^{\mathrm{bin}}],
\qquad
G_6^{\mathrm{bin}}:=\Aut(\Fold_6^{\mathrm{bin}}).
\]
则
\[
\boxed{\;
\dim_{\CC}Z(A_6)=21,
\qquad
\dim_{\FF_2}\bigl((G_6^{\mathrm{bin}})^{\mathrm{ab}}\bigr)=21.\;}
\]
更具体地：
\begin{enumerate}
  \item 任何保持全部 simple block 逐块可分的交换块账本，至少需要 \(21\) 个独立复坐标；该下界由
  \[
  Z(A_6)\cong \CC^{21}
  \]
  精确达到。
  \item 在纤维独立置换语义下，任何完备承载全部奇偶异常信息的二元异常账本，至少需要 \(21\) 个独立 \(\FF_2\)-坐标；该下界由
  \[
  (G_6^{\mathrm{bin}})^{\mathrm{ab}}\cong (\ZZ/2\ZZ)^{21}
  \]
  精确达到。
\end{enumerate}
因此，window-\(6\) 上最小块分离交换坐标数与最小异常比特数严格相等，二者都恰为 \(21\)。
\end{theorem}
\begin{proof}
由定理 \ref{thm:xi-foldbin-groupoid-morita-k-center-dimension-lb}，
\[
Z(A_6)\cong \CC^{F_8}=\CC^{21},
\]
故
\[
\dim_{\CC}Z(A_6)=21.
\]
记 \(e_w\in Z(A_6)\) 为与各个 Wedderburn 块对应的 \(21\) 个中心原始幂等元。它们两两正交且非零，并满足
\[
\sum_{w\in X_6}e_w=1.
\]
若一个交换块账本仍要把全部 simple block 逐块分开，则这些 \(e_w\) 在该账本中的像必须仍然两两可区分；特别地，它们必须给出 \(21\) 个两两正交的非零幂等元。任何有限维交换半单复代数中，两两正交非零幂等元的个数不超过其复维数，故此类账本的交换坐标数必至少为 \(21\)。而中心 \(Z(A_6)\) 本身恰以 \(21\) 个坐标达到该下界。

另一方面，由定理 \ref{thm:xi-foldbin-gauge-group-order-abel-center-closed}，
\[
\bigl(G_6^{\mathrm{bin}}\bigr)^{\mathrm{ab}}\cong (\ZZ/2\ZZ)^{21},
\]
故
\[
\dim_{\FF_2}\bigl((G_6^{\mathrm{bin}})^{\mathrm{ab}}\bigr)=21.
\]
再由定理 \ref{thm:xi-window6-oracle-anomaly-coordinate-bifurcation}，在保持纤维独立置换语义不变的前提下，完备异常信息的最小二元承载对象恰需 \(21\) 个独立 \(\FF_2\)-坐标，任何更小的方案都必然破坏该逐纤维直积语义。

故块分离交换坐标数与异常比特数在 window-\(6\) 上严格对齐，并且都以 \(21\) 为最小值。证毕。
\end{proof}

\begin{proposition}[bin-fold 群胚代数的首次非交换阈值]\label{prop:xi-time-part60acb-binfold-first-noncommutative-threshold}
\leanverified{paper\_xi\_time\_part60acb\_binfold\_first\_noncommutative\_threshold}
对每个 \(m\ge 1\)，记
\[
A_m\cong \bigoplus_{w\in X_m}M_{d_w}(\CC),
\qquad
|X_m|=F_{m+2}.
\]
则：
\begin{enumerate}
  \item 对所有 \(m\ge 1\)，\(A_m\) 都不是单代数；
  \item \(A_m\) 可交换当且仅当 \(m=1\)；
  \item 因而 bin-fold 群胚代数的首次非交换阈值恰发生在
  \[
  \boxed{\;m=2.\;}
  \]
\end{enumerate}
\end{proposition}
\begin{proof}
由 \(F_{m+2}\ge 2\)（\(m\ge 1\)）可知
\[
|X_m|=F_{m+2}\ge 2.
\]
因此 \(A_m\) 的 Wedderburn 分解中至少含有两个简单块，从而 \(A_m\) 对所有 \(m\ge 1\) 都不是单代数。

再看交换性。有限维半单复代数
\[
\bigoplus_{w\in X_m}M_{d_w}(\CC)
\]
可交换，当且仅当每个块都为 \(1\times 1\)，亦即
\[
d_w=1
\qquad (\forall\,w\in X_m).
\]
这等价于
\[
\sum_{w\in X_m}d_w=|X_m|.
\]
而左端恒等于 \(2^m\)，右端恒等于 \(F_{m+2}\)，故交换性等价于
\[
2^m=F_{m+2}.
\]
当 \(m=1\) 时，
\[
2^1=2=F_3,
\]
于是 \(A_1\cong \CC\oplus \CC\) 确为交换。
当 \(m\ge 2\) 时，Fibonacci 数满足
\[
F_{m+2}<2^m,
\]
例如可由归纳得到：\(F_4=3<4=2^2\)，且若
\[
F_{n+2}<2^n,\qquad F_{n+1}<2^{n-1},
\]
则
\[
F_{n+3}=F_{n+2}+F_{n+1}<2^n+2^{n-1}<2^{n+1}.
\]
从而平均块大小
\[
\frac{1}{|X_m|}\sum_{w\in X_m}d_w=\frac{2^m}{F_{m+2}}>1.
\]
故必存在某个 \(w\) 使 \(d_w\ge 2\)，于是 \(A_m\) 非交换。

综上，\(A_m\) 可交换当且仅当 \(m=1\)，故首次非交换阈值恰为 \(m=2\)。证毕。
\end{proof}

\endinput
